
En
effet, on prend @@M0@@ avec @@M1@@, et @@M2@@, de sorte qu'on a une suite exacte
@@M3@@
en vertu de 1.11 et 1.3. De m\^eme, tenant compte de 6.2, on prend @@M4@@ avec @@M5@@, et @@M6@@.

\textbf{Th\'eor\`eme 6.4.} \emph{Sous les conditions de 6.1, on suppose @@M7@@ un quotient de @@M8@@ avec @@M9@@. Soit @@M10@@ pour @@M11@@, o\`u @@M12@@ est le @@M13@@-i\`eme @@M14@@-polyn\^ome (1.9)\,; soit @@M15@@, et soit @@M16@@. Alors}~:
\begin{enumerate}[label=(\roman*)]
\item \emph{@@M17@@ est @@M18@@-r\'egulier.}
\item \emph{@@M19@@}
\item \emph{@@M20@@ est un @@M21@@-faisceau.}
\end{enumerate}

En effet, prenons une section @@M22@@ de @@M23@@ telle qu'on obtienne un diagramme commutatif avec des lignes et colonnes exactes
@@M24@@
